\chapter[Introduction]{Introduction}
\label{Chap:Intro}

% ***************************************************
% Introduction
% ***************************************************

\section{Galaxy Clusters}
\begin{itemize}
    \item What is a Galaxy Cluster (grav bound galaxies, in a bunch of hot gas)
    \item Generic information about galaxy clusters (yay statistics)
    \item Fun facts about galaxy clusters
\end{itemize}

\section{Observational work with Galaxy Clusters}
\begin{itemize}
    \item Well, we can't just go there so we get to look at them really closely.
    \item How do we get data about clusters
    \item In increasing wave length, Optical, X-ray, Radio, Gamma. 
    \item Some telescope + research that we can do with each of these.
\end{itemize}

\section{Radio Emission}
\begin{itemize}
    \item This is what we are interested in, can gain information about clusters from this.
    \item Link synchrotron emission to the electron / gas properies of the cluster and how radio emission can inform us of what the cluster is doing.
    \item Common sources of relativistic electrons are the AGNs, lose energy over time. 
    \item Structures larger than the lifetime of these jets, more diffuse (Halos).
\end{itemize}

\subsection{Giant Radio Halos}
\begin{itemize}
    \item These are of similar extent to the cluster gas. 
    \item They 
\end{itemize}

\subsection{Minihalos}
'\begin{itemize}
    \item What is a minihalo (diffuse radio emission in the core of a cluster.)
    \item Actually they are somewhat common. 
    \item Talk about some previous works on minihalos
    \item Difficulty in distinguishing between production mechanisms.
\end{itemize}

\section{Our goals}
\begin{itemize}
    \item We want to be able to measure minihalo flux for this subset of minihalos at high radio frequencies so that we can investigate if steepening is occuring.
\end{itemize}

\chapter[Observations]{Observations}
This chapter would be about radio interferometry, how it works to some extent and the data that is retrieved from it. How to calibrate the data etc.

\section{What is an interferometer}
\begin{itemize}
    \item Limits of dish size in resolution, causes us to explore different ideas.
    \item How does an interferometer work.
    \item Bunch of dishes being interfered, delay to tune positioning, uv plane to determine coverage.
    \item From all of this we retrieve amp + phase measurements.
\end{itemize}

\section{Our observations}
\begin{itemize}
    \item We are using VLA data in X-band, what does this mean and what are the specs for the VLA.
    \item These observations are of list minihalos and this setup, observation proposal.
    \item Original source of these minihalos is the G14 paper (among others).
    \item Table of minihalos + observing setup.
\end{itemize}


\chapter[Calibration]{Calibration}
This would be how to calibrate and the calibration results for the datasets, plus any special considerations we took.

\section{How to calibrate a dataset}
\begin{itemize}
    \item We have the VLA calibration pipeline to remove stress from our lives.
    \item We have a non-neglible amount of RFI to add stress to our lives. 
    \item A variety of steps but first pain, what is RFI and what does it cause.
\end{itemize}

\section{Flagging data}
\begin{itemize}
    \item Initial testing showed that the pipeline flagging steps didn't work well for our data.
    \item AOFLAGGER worked much better and has a mostly out-of-box VLA script.
    \item How does AOFLAGGER work
    \item Integrating into a pipeline step to still use the pipeline
\end{itemize}

\section{Calibration pipeline}
\begin{itemize}
    \item Work through each calibration step explaining what it does and why it is important. 
    \item Final preparation for imaging is averaging down to further reduce noise.
    \item Notes about general trends from the calibration, with RFI in 10-12 GHz D config being bad, much of it being flagged.
\end{itemize}

\section{Special mentions}
\begin{itemize}
    \item Either here or during following imaging section talk about the 2-3 datasets we processed slightly differently due to bad fits.
    \item Describe the fitting process for the custom bandpass.
\end{itemize}

\chapter[Analysis Tools]{Analysis Tools}
This chapter starts with an overview on imaging and how that is performed. Then we talk a bit about how we actually want to measure the flux of the minihalos.
Could combine with actual results.

\section{Imaging}
\begin{itemize}
    \item When we make a raw image from the data, we still have a lot of artefacts from the point spread function, which comes from the beam.
    \item This means we need a more advanced imaging process to remove this to get a clean image. Implemented in a few different ways, we choose tclean in CASA.
    \item Tclean works by modelling the PSF, creating an empty model, find brightest point in residuals, model that as real, subtract model from residuals, repeat 20K times.
    \item This works mostly well, but we can improve results using self calibration.
\end{itemize}

\subsection{Self calibration}
\begin{itemize}
    \item The phase calibration is the most iffy part for the main calibration, so we can make a test image and check the model seems reasonable.
    \item Then use that model to make small phase corrections to our dataset. Then repeat that a few times and these corrections should converge since it is circular.
    \item This is the final calibration steps that we do, phase corrections 20 degrees $\rightarrow$ 1 degree after a few rounds. 
\end{itemize}

\subsection{Final cleaning}
\begin{itemize}
    \item Between self calibrations this we also check the datasets and do some final manual cleaning.
    \item Clipping the datasets to a sensible amp range, removing spikes, SPWs where there is small amounts of data left.
    \item Overview of the \texttt{assess\_spw} function used for this.
\end{itemize}

\section{Final datasets}
\begin{itemize}
    \item Could have a table here talking about each minihalo, how much data is remaining, spws that are flagged out, some sort of visual. 
    \item Maybe bar graph of spws, line or point graph for minihalo C/D (either total or differentiated by style) normalised to [0, 1]
\end{itemize}

\section{Calibration Error}
\begin{itemize}
    \item Calibration error exists, and sources can vary in their flux amount.
    \item We can try and constrain the error and variance by comparing the secondary calibrator models for D config and the point source fluxes for the images.
    \item Fluxes for secondaries can vary a lot over a year (100\%?) maybe, but lower per month. Can then compare to differences in sources in the image.
    \item If a trend is observed, likely calibration error. If no trend or close models, constrain calibration error and can see the variance in AGN etc.
    \item Can retrieve the point source flux by cutting out diffuse sources (uvrange) and by masking to a tight border.
\end{itemize}

\section{Calculating the minihalo flux}
\begin{itemize}
    \item We can do this in a variety of ways, but we don't want the AGN contamination so we should subtract this out.
    \item Do this per dataset to account for variation over time, incorrect subtractions can result in negative bowls in emission, or remaining flux.
    \item Other problematic sources are dealt with the same way.
    \item Image the subtracted data and then we create a contour for significant emission (3 sigma) from the residual RMS.
    \item Calculate total flux in the region.
    \item Then we calculate the error using this formula from G14 modified for the AGN error since we think its lower.
\end{itemize}

\chapter[Results]{Results}
\begin{itemize}
    \item Sections for each minihalo with a brief explanation of imaging process, custom things we had to do to get results.
    \item Number of self calibration, sources subtracted, difficulties in imaging certain aspects.
    \item Discussion on the calibration error and variance in the observed point sources.
    \item Includes an image of the uvrange data showing point sources, image of the minihalo with contours, and anything else we desire.
    \item Includes discussion compared to previous results or measured values (likely G14) and any caveats or concerns.
\end{itemize}